problem:  
**Conjecture (Homogeneous metrics conformal to Ricci‑flat Kähler metrics are flat).**  
Let \((M^6,g,J)\) be a connected, locally homogeneous Hermitian manifold with integrable complex structure \(J\).  
Assume that there exists a smooth, nonconstant function \(f\) such that the conformally related metric  
\[
\tilde g = e^{2f} g
\]
is Kähler and Ricci‑flat; equivalently, the form \(\tilde \omega = e^{2f}\omega\) with \(\omega(\cdot,\cdot)= g(J\cdot,\cdot)\) satisfies  
\[
d\tilde\omega = 0, \qquad \mathrm{Ric}(\tilde g) = 0.
\]
  
**Conjecture.**  
If \((M,g,J)\) is locally homogeneous, then such a conformal rescaling can exist only if both \(g\) and \(\tilde g\) are flat.  

That is,  
> there are **no nonflat, locally homogeneous Hermitian metrics** whose conformal class contains a Ricci‑flat Kähler (hence Calabi–Yau) metric.

---

Why is it a "good" problem:  

1. **Concise and self‑contained formulation.**  
   The conjecture relies solely on well‑defined objects—locally homogeneous Hermitian metrics and conformal rescaling to a Ricci‑flat Kähler metric. It avoids imprecise torsion language and declares the exact curvature condition (\(\mathrm{Ric}(\tilde g)=0\)) that drives the holonomy reduction to \(SU(3)\).

2. **True conceptual tension.**  
   Local homogeneity implies all curvature data are constant under parallel transport from the standpoint of \(g\), while Ricci‑flat Kähler metrics impose highly nonlinear analytic constraints on the conformal factor \(f\). The conjecture tests whether these two regimes—algebraic symmetry and analytic special holonomy—can intersect nontrivially.  

3. **Bridges disparate theories.**  
   It sits precisely at the interface of  
   - homogeneous Riemannian geometry (Lie‑algebraic curvature),  
   - complex/Kähler geometry (integrability and Ricci‑flatness), and  
   - conformal geometry (transformation of curvature under scaling).  
   Understanding how the homogeneous Killing algebra constrains the conformal factor’s Hessian (entering the conformal Ricci formula) could synthesize these domains in new ways.

4. **Genuinely open and subtle.**  
   While known theorems rule out *homogeneous Ricci‑flat* metrics (which are automatically flat), the conformal variant is far less understood because the group preserving \(g\) need not preserve \(f\). Determining whether any such \(f\) can exist under the curvature constraints remains unclassified and plausibly open.  

5. **Potential unifying implications.**  
   A proof would reinforce the emerging principle that **local homogeneity forbids special holonomy even up to conformal equivalence**, paralleling rigidity conjectures at \(G_2\) and \(\mathrm{Spin}(7)\) levels. A counterexample would have striking consequences: producing the first known instance of a nonflat locally homogeneous geometry whose conformal transform attains Calabi–Yau holonomy—deeply impacting both differential geometry and string‑theoretic compactification models.  

6. **Simplicity concealing analytical depth.**  
   The statement fits in one line: *“Can a locally homogeneous metric be conformally Ricci‑flat Kähler and nonflat?”* Yet resolving it demands analyzing the coupled PDEs for \(f\)
   \[
   \nabla^2 f + df\otimes df - \tfrac{1}{6} g|\nabla f|^2 = 0,
   \]
   adapted to a homogeneous context and constrained by the Lie algebraic curvature equations.  
   This blend of precision, brevity, and extensive cross‑field reach captures the essence of a “narrow‑entrance, profound‑depth” research-level mathematical problem.